\documentclass[11pt]{article}

\usepackage[margin=0.85in]{geometry}
\usepackage{amsmath,amssymb}
\usepackage{enumitem}
\usepackage{xcolor}
\usepackage{tcolorbox}
\usepackage{fancyhdr}
\usepackage{lastpage}
\usepackage{hyperref}
\usepackage{microtype}

\definecolor{uwpurple}{HTML}{4B2E83}
\definecolor{uwgold}{HTML}{B7A57A}
\definecolor{softpurple}{HTML}{F4F1FA}

\hypersetup{colorlinks=true, linkcolor=uwpurple, urlcolor=uwpurple}

\pagestyle{fancy}
\fancyhf{}
\lhead{\textbf{MATH 394: Probability I}}
\rhead{\textbf{Homework 8}}
\cfoot{\thepage\ of \pageref{LastPage}}

\setlength{\parindent}{0pt}
\setlength{\parskip}{0.45em}

\newcommand{\course}{MATH 394: Probability I}
\newcommand{\instructor}{Arman Jahangiri}
\newcommand{\department}{Department of Mathematics}
\newcommand{\university}{University of Washington}

\newtcolorbox{problemheader}{
    colback=softpurple,
    colframe=uwpurple,
    boxrule=0.8pt,
    arc=2mm,
    left=2mm,
    right=2mm,
    top=1mm,
    bottom=1mm
}

\newcommand{\source}[2]{%
    \vspace{0.15cm}
    {\small\textit{Source: Blitzstein and Hwang, \textit{Introduction to Probability}, Chapter #1, Problem #2.}}
}

\newcommand{\answerbox}{%
\begin{flushright}
\begin{tcolorbox}[width=0.36\textwidth,height=2cm,colback=white,colframe=uwpurple,boxrule=0.7pt,arc=1mm]
{\small\textbf{Final Answer}}
\end{tcolorbox}
\end{flushright}
}

\newcommand{\partanswerbox}[1]{%
\begin{flushright}
\begin{tcolorbox}[width=0.36\textwidth,height=2cm,colback=white,colframe=uwpurple,boxrule=0.7pt,arc=1mm]
{\small\textbf{Final Answer for Part #1}}
\end{tcolorbox}
\end{flushright}
}

\begin{document}

\begin{titlepage}
\centering
\vspace*{1.2cm}

{\Large \textbf{\university}}\\
{\large \department}

\vspace{1.5cm}

{\Huge \color{uwpurple}\textbf{Homework 8}}\\[0.25cm]
{\Large \textbf{\course}}\\[0.25cm]
{\large Instructor: Arman Jahangiri}

\vspace{1cm}

\begin{tcolorbox}[width=0.78\textwidth,colback=softpurple,colframe=uwpurple,boxrule=0.9pt,arc=3mm]
\centering
\textbf{Submission:} Single PDF on Gradescope

\textbf{Deadline:} 10:00 PM Pacific Time on the listed due date
\end{tcolorbox}

\vspace{0.8cm}

\begin{tcolorbox}[width=0.86\textwidth,colback=white,colframe=uwgold,boxrule=0.8pt,arc=3mm]
This homework is worth \textbf{50 points}.

Across the quarter, there are \textbf{eight homework assignments}, worth \textbf{400 points total}. Homework assignments together account for \textbf{40\%} of the final course grade. Further course policies can be seen in the \href{https://armanjg.github.io/MATH-394-Probability-1-v2/}{\textbf{MATH 394 syllabus}.}
\end{tcolorbox}

\vfill
\end{titlepage}

\section*{Homework Policy}

\subsection*{Submission and Deadline}

Please:

\begin{itemize}[leftmargin=*]
    \item Submit homework as a single PDF on Gradescope by \textbf{10:00 PM (Pacific Time)} on the listed due date.
    \item Begin each problem on a new page, and ensure that pages are correctly tagged in Gradescope.
    \item Write solutions clearly and legibly.
    \item Typesetting solutions in LaTeX is encouraged, but not required.
\end{itemize}

\subsection*{Late Days}

Each student is allotted \textbf{six late days} for the quarter. A late day extends a homework deadline by up to 24 hours without penalty. For example, submitting an assignment anytime between 10:01 PM on the due date and 10:00 PM the following day counts as one late day.

The following rules apply:
\begin{itemize}[leftmargin=*]
   \item No explanation is required to use late days.
   \item At most \textbf{two late days} may be used on any single homework assignment.
   \item Late days may not be used on the final homework assignment or on any assignment for which solutions have already been released.
   \item It is the student's responsibility to keep track of how many late days have been used during the quarter.
\end{itemize}

Once all late days have been exhausted, additional late submissions will incur a penalty of \textbf{10\% per day}, up to a maximum deduction of \textbf{50\%}. Assignments submitted more than five days late, or after solutions have been released, will not be accepted.

\subsection*{Submission Issues and Technical Difficulties}

If a serious technical issue prevents a timely Gradescope submission, students may temporarily submit their assignment by email to the instructor at \href{mailto:armanjg@uw.edu}{armanjg@uw.edu}. In such cases:

\begin{itemize}[leftmargin=*]
    \item The submission must consist of a \textbf{single PDF file}. The email subject line must follow the format:
    \[
    \texttt{HW\# Submission - FULL NAME - STUDENT ID}.
    \]
    \item Once the Gradescope issue has been resolved, the student must upload the \emph{exact same file} to Gradescope as soon as possible.
    \item The student should clearly indicate in the Gradescope submission that the assignment was previously submitted by email before the deadline.
\end{itemize}

Email submissions are intended only for genuine technical emergencies and should not be used as a substitute for timely Gradescope submission.

\newpage

% ------------------------------------------------------------
 

% ------------------------------------------------------------

\begin{problemheader}
\textbf{Problem 1.}
\end{problemheader}
\source{8}{1--5}

Find the following PDFs.

\begin{enumerate}[label=(\alph*), leftmargin=*]
    \item Find the PDF of $e^{-X}$ for $X\sim \operatorname{Exp}(1)$.

    \item Find the PDF of $X^7$ for $X\sim \operatorname{Exp}(\lambda)$.

    \item Find the PDF of $Z^3$ for $Z\sim N(0,1)$.

    \item Find the PDF of $Z^4$ for $Z\sim N(0,1)$.

    \item Find the PDF of $|Z|$ for $Z\sim N(0,1)$.
\end{enumerate}
 
\newpage

% ------------------------------------------------------------

\begin{problemheader}
\textbf{Problem 2. }
\end{problemheader}
\source{8}{13}

Let $X$ and $Y$ be i.i.d. $\operatorname{Exp}(\lambda)$, and let
\[
T=\log(X/Y).
\]
Find the CDF and PDF of $T$.
 
\newpage

% ------------------------------------------------------------

\begin{problemheader}
\textbf{Problem 3.  }
\end{problemheader}
\source{8}{14}

Let $X$ and $Y$ have joint PDF $f_{X,Y}(x,y)$, and transform $(X,Y)\mapsto (T,W)$ linearly by letting
\[
T=aX+bY,
\qquad
W=cX+dY,
\]
where $a,b,c,d$ are constants such that
\[
ad-bc\ne0.
\]

\begin{enumerate}[label=(\alph*), leftmargin=*]
    \item Find the joint PDF $f_{T,W}(t,w)$ in terms of $f_{X,Y}$, though your answer should be written as a function of $t$ and $w$.

    \item For the case where
    \[
    T=X+Y,
    \qquad
    W=X-Y,
    \]
    show that
    \[
    f_{T,W}(t,w)
    =
    \frac12
    f_{X,Y}\left(\frac{t+w}{2},\frac{t-w}{2}\right).
    \]
\end{enumerate}
 
\newpage

% ------------------------------------------------------------

% ------------------------------------------------------------

\begin{problemheader}
\textbf{Problem 4. }
\end{problemheader}
\source{10}{2}

For i.i.d. random variables $X_1,\dots,X_n$ with mean $\mu$ and variance $\sigma^2$, give a value of $n$, as a specific number, that will ensure that there is at least a $99\%$ chance that the sample mean will be within $2$ standard deviations of the true mean $\mu$.
 
\newpage

% ------------------------------------------------------------

\begin{problemheader}
\textbf{Problem 5.  }
\end{problemheader}
\source{10}{4}

The famous arithmetic mean--geometric mean inequality says that for any positive numbers
\[
a_1,a_2,\dots,a_n,
\]
we have
\[
\frac{a_1+a_2+\cdots+a_n}{n}
\ge
(a_1a_2\cdots a_n)^{1/n}.
\]

Show that this inequality follows from Jensen's inequality, by considering $E(\log X)$ for a random variable $X$ whose possible values are $a_1,\dots,a_n$. You should specify the PMF of $X$.
 
\newpage

% ------------------------------------------------------------

\begin{problemheader}
\textbf{Problem 6.  }
\end{problemheader}
\source{10}{7}

Let $X$ and $Y$ be i.i.d. positive random variables, and let $c>0$. For each part below, fill in the appropriate equality or inequality symbol.

Write $=$ if the two sides are always equal, $\le$ if the left-hand side is less than or equal to the right-hand side but not necessarily equal, and similarly for $\ge$. If no relation holds in general, write $?$.

\begin{enumerate}[label=(\alph*), leftmargin=*]
    \item $E(\log X)\ \underline{\hspace{1cm}}\ \log(E X)$
    \item $E(X)\ \underline{\hspace{1cm}}\ \sqrt{E(X^2)}$
    \item $E(\sin^2 X)+E(\cos^2 X)\ \underline{\hspace{1cm}}\ 1$
    \item $E(|X|)\ \underline{\hspace{1cm}}\ \sqrt{E(X^2)}$
    \item $P(X>c)\ \underline{\hspace{1cm}}\ \dfrac{E(X^3)}{c^3}$
    \item $P(X\le Y)\ \underline{\hspace{1cm}}\ P(X\ge Y)$
    \item $E(XY)\ \underline{\hspace{1cm}}\ \sqrt{E(X^2)E(Y^2)}$
    \item $P(X+Y>10)\ \underline{\hspace{1cm}}\ P(X>5\text{ or }Y>5)$
    \item $E(\min(X,Y))\ \underline{\hspace{1cm}}\ \min(EX,EY)$
    \item $E(X/Y)\ \underline{\hspace{1cm}}\ \dfrac{EX}{EY}$
    \item $E\!\left(X^2(X^2+1)\right)\ \underline{\hspace{1cm}}\ E\!\left(X^2(Y^2+1)\right)$
    \item $E\!\left(\dfrac{X^3}{X^3+Y^3}\right)\ \underline{\hspace{1cm}}\ E\!\left(\dfrac{Y^3}{X^3+Y^3}\right)$
\end{enumerate}
 
\newpage

% ------------------------------------------------------------
 

% ------------------------------------------------------------

\begin{problemheader}
\textbf{Problem 7.}
\end{problemheader}
\source{10}{22}

Let
\[
U_1,U_2,\dots,U_{60}
\]
be i.i.d. $\operatorname{Unif}(0,1)$ and let
\[
X=U_1+U_2+\cdots+U_{60}.
\]

\begin{enumerate}[label=(\alph*), leftmargin=*]
    \item Which important distribution is the distribution of $X$ very close to? Specify what the parameters are, and state which theorem justifies your choice.

    \item Give a simple but accurate approximation for
    \[
    P(X>17).
    \]
    Justify briefly.
\end{enumerate}
 
\newpage

% ------------------------------------------------------------

\begin{problemheader}
\textbf{Problem 8.}
\end{problemheader}
\source{10}{25}

\begin{enumerate}[label=(\alph*), leftmargin=*]
    \item Let
    \[
    Y=e^X,
    \qquad
    X\sim \operatorname{Exp}(3).
    \]
    Find the mean and variance of $Y$.

    \item For $Y_1,\dots,Y_n$ i.i.d. with the same distribution as $Y$ from part (a), what is the approximate distribution of the sample mean
    \[
    \bar Y_n=\frac1n\sum_{j=1}^{n}Y_j
    \]
    when $n$ is large?
\end{enumerate}
 
\newpage

% ------------------------------------------------------------

\begin{problemheader}
\textbf{Problem 9. }
\end{problemheader}
\source{10}{27}

Consider i.i.d. $\operatorname{Pois}(\lambda)$ random variables
\[
X_1,X_2,\dots.
\]
The MGF of $X_j$ is
\[
M(t)=e^{\lambda(e^t-1)}.
\]

\begin{enumerate}[label=(\alph*), leftmargin=*]
    \item Find the MGF $M_n(t)$ of the sample mean
    \[
    \bar X_n=\frac1n\sum_{j=1}^{n}X_j.
    \]

    \item Find the limit of $M_n(t)$ as $n\to\infty$. You can do this with almost no calculation using a relevant theorem, or you can use part (a) and the fact that $e^x\approx 1+x$ if $x$ is very small.
\end{enumerate}
 

\end{document}